\documentclass[10pt,a4paper]{amsart}
\usepackage[margin=19mm]{geometry}
\usepackage{amsmath,amssymb,mathtools}
\usepackage[scr=rsfs,cal=euler]{mathalfa}
\usepackage{xfrac}
\usepackage[unicode,hidelinks,bookmarks=false]{hyperref}
\newcommand{\ie}{i.e.\ }
\newcommand{\cf}{cf.\ }
\newcommand{\resp}{respectively\ }
\newcommand{\Subsection}{\subsection}
\providecommand{\nobreakdash}{\nobreak\hbox{-}}
\setlength{\emergencystretch}{2em}
\multlinegap0pt
\usepackage{bm}
\usepackage{bbm}
\usepackage{dsfont}
\usepackage{xspace}
\usepackage{cancel}
\usepackage{mathtools}
\usepackage[shortlabels]{enumitem}
\usepackage{amsthm}
\makeatletter
\@ifundefined{thm}{\newtheorem{thm}{Thm}}{}
\@ifundefined{lemma}{\newtheorem{lemma}[thm]{Lemma}}{}
\makeatother
\DeclareMathOperator{\Int}{Int}
\def\Z{{\mathbb{Z}}}
\def\cV{\mathcal {V}}
\def\cW{\mathcal {W}}
\newcommand{\clo}{{\text{cl}}\,\,}
\newcommand{\sm}{\setminus}
\renewcommand{\tilde}{\widetilde}
\title[Transcendental correspondences: when Fuchsian groups take ov]{Transcendental correspondences: when Fuchsian groups take over basins of entire maps}
\author{Kostiantyn Drach, Sabyasachi Mukherjee, Leticia Pardo-Sim\'on}
\date{Source: arXiv:2607.06090v1 (CC BY 4.0)}
\begin{document}
\maketitle

\noindent\emph{Source and licence.} Transcendental correspondences: when Fuchsian groups take over basins of entire maps. Version arXiv:2607.06090v1 (CC BY 4.0), distributed under the Creative Commons Attribution 4.0 International licence, as recorded in the arXiv licence metadata for this exact version and shown on the abstract page https://arxiv.org/abs/2607.06090v1. Full rights notice: licenses/arxiv-2607.06090-CC-BY-4.0.txt.\par\smallskip

\noindent\textit{Editorial scope.} Counted unit: the Lemma whose source label is \texttt{lifted\_tiling\_top\_lem} in the article (source0.tex lines 1700--1707, source label \texttt{lifted\_tiling\_top\_lem}), delivered together with the article's own complete original proof of it (source0.tex lines 1708--1727, one \texttt{proof} environment of 20 lines). No mathematics is written, repaired or re-derived: statement and proof are the article's own text line for line. Required context retained uncounted: none. Every definition, display and label of those lines is retained unchanged. Omitted from this package and not redistributed: 1--1699, 1728--2087. Nothing inside a retained excerpt is omitted or altered: every retained body is delivered exactly as the article's lines are written, so no omission receipt of any kind accompanies this package. Citations: the retained excerpts carry no citation command at all, so no bibliography entry of the article is transcribed and no reference is redistributed. Reference adaptation: none was needed and none was made; every reference inside the delivered unit resolves inside it, so no unresolved reference or citation is produced. Printed slips kept literal: no printed word, symbol, hypothesis or number of the counted statement or of its proof was altered. Layout and class: the article's own class is \texttt{amsart} with packages including xcolor, braket, comment, soul, hyperref, pdfpages, faktor, graphicx, mathabx, bbm, url, caption, subcaption, mathrsfs; the wrapper is the lane's pinned \texttt{amsart} editorial wrapper at 10pt on A4, which restates the theorem-like environments the retained text needs and adds the article's own macro definitions that the retained text needs (\texttt{\textbackslash Int}, \texttt{\textbackslash Z}, \texttt{\textbackslash cV}, \texttt{\textbackslash cW}, \texttt{\textbackslash clo}, \texttt{\textbackslash sm}, \texttt{\textbackslash tilde}), with extra wrapper packages (bm, bbm, dsfont, xspace, cancel, mathtools, enumitem). The delivered numbering is the unit's own. Assets: no retained span contains a figure environment, a graphics-inclusion command, a PDF-page inclusion, a table or a TikZ picture; no image file is copied into this package, the package declares no asset dependency and no image file is redistributed with it. Licence: version arXiv:2607.06090v1 of this article is distributed under the Creative Commons Attribution 4.0 International licence, as recorded in the arXiv licence metadata for this exact version and shown on the abstract page https://arxiv.org/abs/2607.06090v1. A full rights notice is delivered at licenses/arxiv-2607.06090-CC-BY-4.0.txt. Characters: every retained excerpt is pure ASCII, so the delivered engine, XeLaTeX with its default Latin Modern text font, reports no missing character.
\begin{lemma}\label{lifted_tiling_top_lem}
\noindent\begin{enumerate}[leftmargin=8mm]
\item The set $\displaystyle\clo{\pmb{\cW}}=\bigcup_{i=0}^{p-1} \overline{\pmb{\cW}_i}$ is connected, where each $\pmb{\cW}_i$ is a simply connected domain (in fact, a Jordan domain if $p\neq 1$) that is mapped as a degree $n=n(\Sigma)$ branched covering onto $\tilde U$ by $\mathfrak{h}$.
\item  If $p\neq 1$, then $\overline{\pmb{\cW}_i}\cap \overline{\pmb{\cW}}_{i+1}$ is a single point.
\item If $p\neq 1$, then $\overline{\pmb{\cW}_i}\cap \overline{\pmb{\cW}_j}=\emptyset$ if $\vert j-i\vert\neq1$, $i,j\in\Z/p\Z$.
\item If $p=1$, then the boundary $\pmb{\cW}$ is topologically a figure-eight curve.
\end{enumerate}
\end{lemma}

\begin{proof}
It follows from the construction of the conformal combination $g$ that $\clo\tilde U\,\setminus\, \tilde T$ is a connected set. In fact, it is the union of $p$ closed topological disks $Y_0,\cdots,Y_{p-1}$ (except in the case $p=1$, when $Y_0=\clo\tilde U\,\setminus\, \tilde T$ is the closure of a simply connected domain such that there is a unique boundary point with two accesses from the domain and all other boundary points have a unique access from the domain)  such that, after possibly renumbering, we have the following properties:
\begin{itemize}[leftmargin=8mm]
\item $\partial Y_i\cap\partial \tilde T$ is a non-singular real-analytic curve that is a Jordan curve for $p=1$ and an arc for $p\neq 1$,
\item $Y_i\cap Y_{i+1}$ is a single point on $\partial\tilde U$, and
\item $Y_i\cap Y_j=\emptyset$ if $\vert j-i\vert\neq1$, where $i,j\in\Z/p\Z$.
\end{itemize}
The inverse branch $(\mathfrak{h}\vert_{\overline{\cV^0}})^{-1}$ maps each $Y_i$ to a closed topological disk $\pmb{Y}_i\subset \overline{\cV^0}$ (except for $p=1$, when $\pmb{Y}_0$ is the closure of a simply connected domain such that there is a unique boundary point with two accesses from the domain and all other boundary points have a unique access from the domain). By construction, the boundary $\partial\cV^0$ is contained in the boundary of $\left(\mathfrak{h}\vert_{\overline{\cV^0}}\right)^{-1}(\clo \tilde U\,\sm\,\tilde T)=\bigcup_{i=1}^p \pmb{Y}_i$. More precisely, $\bigcup_{i=1}^p \pmb{Y}_i$ is a connected set  containing a relative neighborhood of $\partial\cV^0\setminus\{x_1,\cdots,x_p\}$ in $\overline{\cV^0}$, where $\{x_1,\cdots,x_p\}\subset\partial\cV^0$ is $\eta-$invariant.
By definition of $\pmb{\cW}$, we have that
\begin{equation*}
\clo \pmb{\mathcal{W}} = \Bigg(\big(\mathfrak{h}\big|_{\overline{\cV^0}}\big)^{-1}(\clo\tilde U\,\sm\tilde T)\Bigg)\bigcup\Bigg(\eta\Big(\big(\mathfrak{h}\big|_{\overline{\cV^0}}\big)^{-1}(\clo\tilde U\,\sm\tilde T)\Big)\Bigg)
= \Big(\bigcup_{i=1}^p \pmb{Y}_i\Big)\bigcup\Bigg(\eta\Big(\bigcup_{i=1}^p \pmb{Y}_i\Big)\Bigg).
\end{equation*}
For $i\in\{0,\cdots,p-1\}$, we define $\pmb{\cW}_i$ to be 
$$
\pmb{\cW}_i:=\Int{\Bigg(\overline{\Int{\pmb{Y}_i}\bigsqcup\eta(\Int{\pmb{Y}_{i'}})}\Bigg)},
$$ 
where $i'\in\{0,\cdots,p-1\}$ is the unique index such that $g(\partial Y_i\cap\partial\tilde T)=\partial Y_{i'}\cap\partial\tilde T$.
The desired topological properties of $\pmb{\cW}_i$ now follows from the above description of $\pmb{\cW}_i$ and the $\eta-$invariance of $\{x_1,\cdots,x_p\}$.
\end{proof}

\end{document}
